\section{照料、身份与暴力：大事如何落在不同行动者身上}
\label{sec:synthesis-care-status-violence}

编年史常以皇帝、将军和大城的命运推进，读者却不应因此把家庭、性别、年龄和身份视作附属细节。迁徙时谁照料老人和幼儿，征役后谁维持田作，围城中谁有资格分到粮食，战败后谁会被俘、被卖、被迫改换依附，都会影响一个政权的资源与合法性。材料很少允许我们替无名者写出完整传记；它仍足以使我们拒绝一种假设，即所有人以同样方式经历战争和制度。

\subsection{家内照料是政治选择的条件}

李密《陈情表》把侍奉祖母与接受新朝征辟写成必须请求皇帝裁断的冲突。它不是一份普通家庭的调查，却清楚显示国家任用一名士人要穿过家内照料与伦理语言。类似的限制也存在于没有留下表章的人家：成年人被征发、迁移或死亡时，儿童、老人、病者与留守者的照料需要由亲属、邻里、依附关系或寺院承担。国家账簿所见的是丁口与役使，家庭真正面对的则是失去一个劳力后如何继续生活。

南渡、侨置和土断改变的也包括这种照料网络。旧乡的亲属、坟墓、田地与邻里未必能随人渡江，新住处又未必立刻提供可信赖的关系。婚姻、寄居、收养与共同劳作可能成为重新安置的方式，但材料往往只在墓志、家训或法律条文的边缘留下痕迹。我们不能从个别精英家庭推出普遍经验，却可以指出：若把迁徙只写成男性士人的仕宦选择，便会遗漏使这种选择成为可能或付出代价的人。

\subsection{身份不是天然不变的保护}

两晋南北朝文本中有官人、良民、部众、军户、僧尼、奴婢、降人和流民等多种称谓。它们既可以提供某种资源与保护，也可能在危机中成为征发、排斥或暴力的依据。不能把这些类别直接换成固定的现代社会群体；应当问它在当时的登记、军政、法律或叙事中意味着什么。一个“降人”在战报里或许只是胜者的数量，在实际安置中却可能面对新的土地、劳役和亲属分离。

冉闵时期针对被分类人群的暴力，说明称谓并不只是史家的描述。政治竞争能够把“谁属于何类”变成杀戮与动员的语言，受害范围和死亡数字又难以从传世叙事中精确统计。承认数字与群体界定的争议，绝不是淡化暴力；相反，它防止后人用不稳的总数和本质化标签遮蔽暴力怎样在城市、道路和家庭中发生。受害者并非只属于一张族名表，他们也是逃离、被捕、失去亲属或被迫投靠的人。

\subsection{宗教空间与脆弱者}

寺院在某些时候可以提供施舍、安置、医疗、葬仪或跨地域联系，因此可能成为失去家族支持者的一个节点；这并不意味着它们是人人都能进入的慈善系统。供养能力、僧团规则、地方政治和战乱条件都会限制其作用。禁佛或城市陷落时，寺院人员和资源同样会被清查、驱散或掠夺。宗教史若只记录帝王信仰，便会错过寺院在照料与脆弱之间所处的实际位置。

女性的宗教供养与墓志记录提示她们并非只在男性叙事的背后。题记和碑文有时保留女性作为施主、母亲、妻子或亡者的名字，显示家族、功德与哀悼可以通过她们进入公共文字。但这些记录仍受书写者与资源的筛选，不能被用来断言女性普遍拥有同等财产或行动自由。让有限证据保持有限，才使可见的行动不被夸张，也使不可见的劳动继续成为应当追问的问题。

\begin{debate}[怎样写没有自述的人]
没有第一人称文字，不等于没有历史行动；为沉默者虚构内心，也不是补救。较可靠的办法是从制度、道路、征发、墓志称谓、题记和物质遗存提出受约束的问题：谁可能承担何种劳作，哪些关系因迁徙或战争被打断，哪些结论仍无法证实。这样既不把无名者消失，也不把想象伪装成证词。
\end{debate}

\subsection{暴力之后的长期时间}

城破、屠杀、俘获和征发的直接后果最容易进入年表，长期后果往往更难追踪。一个家庭失去劳力后，可能改变耕作、婚姻或依附；一座城市失去工匠和商人后，重建所需的时间可能远长于新政权宣布接收的日期；一条道路因不安全而停止通行，又会改变宗教、文本和市场的流动。材料无法为每种后果作完整计算，却要求我们不要把战争结尾写成生活结尾。

因此，照料、身份与暴力应贯穿而非附在本卷的六层叙事之后。它们使制度得失具有真实的承担者：一项登记给谁带来可预期的责任，一次迁都使谁离开支持网络，一场边境战争由谁运输粮食，一次禁令让哪些人失去可依靠的场所。把问题如此落地，国家史才不会只是一连串没有人生活其中的名称。

这也是综合章回到编年的原因。永嘉、淝水、六镇、侯景与陈亡不是只有君主和将领的节点；每一个节点都重新分配了保护、劳役、照料和逃离的可能。我们未必知道每一个人的选择，却能从道路的阻断、户籍的重编、寺院的处置和城市的接收中看出，政治权力如何改变了家庭可以依赖的世界。只有把这些后果留在年序里，才能既不虚构沉默者的声音，也不让他们消失在王朝更替之外。
